% ID: MAT-STA-MMM-001
% Titolo: Statistiche di una serie di dati simulati
% Disciplina: Matematica
% Area: statistica
% Argomento: media_mediana_moda
% Classe: Prima secondaria di secondo grado
% Difficolta: 2
% Tipo: laboratorio
% Tag: matematica, statistica, media_mediana_moda
% Fonte: laboratorio_dati_anonimizzato
% Provenienza: materiale fornito dal docente; testo anonimizzato; soluzione redatta per Exergo

\begin{esercizio}
Sono dati i valori simulati di temperatura $68,61,55,50,46,43$ gradi Celsius,
rilevati rispettivamente ai minuti $0,2,4,6,8,10$.
Calcola media aritmetica, massimo, minimo e variazione totale. Rappresenta le coppie con un grafico XY.
Come cambiano i risultati se l'ultimo valore diventa 40? Ripristina poi 43.
\end{esercizio}

\begin{soluzione}
La somma è 323, quindi media $323/6\simeq53{,}83$, massimo 68, minimo 43,
variazione $43-68=-25$. Con ultimo valore 40: media $320/6\simeq53{,}33$,
massimo 68, minimo 40, variazione $-28$. Il grafico passa per le sei coppie assegnate
ed è decrescente. In Excel, dati in B5:B10: \texttt{=MEDIA(B5:B10)},
\texttt{=MAX(B5:B10)}, \texttt{=MIN(B5:B10)}, \texttt{=B10-B5}.
\end{soluzione}
